\begin{thebibliography}{00}

\bibitem{soni2025}
M. S. Soni, K. Agrawal, P. Parikh, P. Bharadva, and M. Myangar,
``Empowering college girls: From research insights to AI-based solutions
for menstrual health and academic success,''
\textit{Journal of Neonatal Surgery}, vol.~14, no.~32s,
pp.~7896--7902, 2025.

\bibitem{li2021}
K. Li, I. Urteaga, A. Shea, V. J. Vitzthum, C. H. Wiggins, and N. Elhadad,
``A generative, predictive model for menstrual cycle lengths that accounts
for potential self-tracking artifacts in mobile health data,''
arXiv preprint arXiv:2102.12439, 2021.

\bibitem{wang2024}
D. Wang and S. Zhang,
``Large language models in medical and healthcare fields: applications,
advances, and challenges,''
\textit{Artificial Intelligence Review}, vol.~57, p.~299, 2024,
doi: 10.1007/s10462-024-10921-0.

\bibitem{rossi2025}
M. Rossi and S. Rehman,
``Integrating artificial intelligence into telemedicine: Evidence, challenges,
and future directions,''
\textit{Cureus}, vol.~17, no.~8, e90829, 2025,
doi: 10.7759/cureus.90829.

\bibitem{epstein2017}
D. A. Epstein, N. B. Lee, J. H. Kang, E. Agapie, J. Schroeder,
L. R. Pina, J. Fogarty, J. A. Kientz, and S. A. Munson,
``Examining menstrual tracking to inform the design of personal informatics
tools,''
in \textit{Proceedings of the 2017 CHI Conference on Human Factors in
Computing Systems}, 2017, pp.~6876--6888,
doi: 10.1145/3025453.3025635.

\bibitem{rajesh2025}
M. Rajesh,
``Adaptive edge-federated AI framework for contactless menstrual health
prediction using multimodal physiological intelligence,''
\textit{MethodsX}, vol.~15, p.~103665, 2025,
doi: 10.1016/j.mex.2025.103665.

\bibitem{ali2023}
M. M. Ali, H. A. A. Algashamy, E. Alzidi, K. Ahmed,
F. M. Bui, S. K. Patel, S. Azam, L. F. Abdulrazak, and M. A. Moni,
``Development and performance analysis of machine learning methods for
predicting depression among menopausal women,''
\textit{Healthcare Analytics}, vol.~3, p.~100202, 2023,
doi: 10.1016/j.health.2023.100202.

\bibitem{yu2022}
J.-L. Yu, Y.-F. Su, C. Zhang, L. Jin, X.-H. Lin, L.-T. Chen,
H.-F. Huang, and Y.-T. Wu,
``Tracking of menstrual cycles and prediction of the fertile window via
measurements of basal body temperature and heart rate as well as
machine-learning algorithms,''
\textit{Reproductive Biology and Endocrinology}, vol.~20, p.~118, 2022,
doi: 10.1186/s12958-022-00993-4.

\bibitem{iturriago2025}
L. M. Iturriago-Salas, J. A. Mesa-Sarmiento, P. A. Castro-Cabrera,
A. M. \'{A}lvarez-Meza, and G. Castellanos-Dominguez,
``Artificial intelligence-driven mobile platform for thermographic imaging
to support maternal health care,''
\textit{Computers}, vol.~14, p.~466, 2025,
doi: 10.3390/computers14110466.

\bibitem{abbas2025}
Q. Abbas, W. Jeong, and S. W. Lee,
``Explainable AI in clinical decision support systems: A meta-analysis of
methods, applications, and usability challenges,''
\textit{Healthcare}, vol.~13, p.~2154, 2025,
doi: 10.3390/healthcare13172154.

\bibitem{kim2024}
H.-K. Kim,
``The effects of artificial intelligence chatbots on women's health:
A systematic review and meta-analysis,''
\textit{Healthcare}, vol.~12, p.~534, 2024,
doi: 10.3390/healthcare12050534.

\bibitem{kotz2016}
D. Kotz, C. A. Gunter, S. Kumar, and J. P. Weiner,
``Privacy and security in mobile health: A research agenda,''
\textit{Computer}, vol.~49, no.~6, pp.~22--30, 2016,
doi: 10.1109/MC.2016.185.

\bibitem{khairunisa2025}
M. Khairunisa, T. Hasanah, and R. Anwar,
``Comparison of machine learning models for menstrual cycle prediction,''
\textit{Journal of Applied Informatics and Computing}, 2025.
Available: \url{https://jurnal.polibatam.ac.id/index.php/JAIC/article/view/9076}

\bibitem{zad2024}
Z. Zad, V. S. Jiang, A. T. Wolf, T. Wang, J. J. Cheng,
I. C. Paschalidis, and S. Mahalingaiah,
``Predicting polycystic ovary syndrome with machine learning algorithms
from electronic health records,''
\textit{Frontiers in Endocrinology}, vol.~15, p.~1298628, 2024,
doi: 10.3389/fendo.2024.1298628.

\bibitem{pi2025}
X. Pi, J. Wang, L. Chu, G. Zhang, and W. Zhang,
``Prediction of high-risk pregnancy based on machine learning algorithms,''
\textit{Scientific Reports}, vol.~15, p.~15561, 2025,
doi: 10.1038/s41598-025-00450-3.

\bibitem{ahmed2023}
S. Ahmed, M. S. Rahman, I. Jahan, M. S. Kaiser,
A. S. M. S. Hosen, D. Ghimire, and S.-H. Kim,
``A review on the detection techniques of polycystic ovary syndrome
using machine learning,''
\textit{IEEE Access}, vol.~11, pp.~86522--86543, 2023,
doi: 10.1109/ACCESS.2023.3304536.

\bibitem{rahman2024}
M. M. Rahman, A. Islam, F. Islam, M. Zaman,
M. R. Islam, M. S. A. Sakib, and H. M. H. Babu,
``Empowering early detection: A web-based machine learning approach
for PCOS prediction,''
\textit{Informatics in Medicine Unlocked}, vol.~47, p.~101500, 2024,
doi: 10.1016/j.imu.2024.101500.

\bibitem{maity2025}
S. Maity and M. J. Saikia,
``Large language models in healthcare and medical applications: A review,''
\textit{Bioengineering}, vol.~12, p.~631, 2025,
doi: 10.3390/bioengineering12060631.

\bibitem{lundberg2017}
S. M. Lundberg and S.-I. Lee,
``A unified approach to interpreting model predictions,''
in \textit{Advances in Neural Information Processing Systems 30},
2017, arXiv:1705.07874.

\bibitem{goodale2019}
B. M. Goodale, M. Shilaih, L. Falco, F. Dammeier,
G. Hamvas, and B. Leeners,
``Wearable sensors reveal menses-driven changes in physiology
and enable prediction of the fertile window: Observational study,''
\textit{Journal of Medical Internet Research},
vol.~21, no.~4, p.~e13404, 2019,
doi: 10.2196/13404.

\bibitem{urteaga2017}
I. Urteaga, D. J. Albers, M. V. Wheeler, A. Druet,
H. Raffauf, and N. Elhadad,
``Towards personalized modeling of the female hormonal cycle:
Experiments with mechanistic models and Gaussian processes,''
in \textit{Advances in Neural Information Processing Systems}, 2017,
arXiv:1712.00117.

\bibitem{liu2019}
B. Liu, S. Shi, Y. Wu, D. Thomas, L. Symul,
E. Pierson, and J. Leskovec,
``Predicting pregnancy using large-scale data from a women's
health tracking mobile application,''
in \textit{Proceedings of the 2019 World Wide Web Conference
(WWW '19)}, 2019, pp.~1--7,
doi: 10.1145/3308558.3313512.

\bibitem{kilungeja2025}
G. Kilungeja, K. Graham, X. Liu, and M. Nasseri,
``Machine learning-based menstrual phase identification using
wearable device data,''
\textit{npj Women's Health}, vol.~3, p.~29, 2025,
doi: 10.1038/s44294-025-00078-8.

\bibitem{li2020}
K. Li, I. Urteaga, C. H. Wiggins, A. Druet, A. Shea,
V. J. Vitzthum, and N. Elhadad,
``Characterizing physiological and symptomatic variation in
menstrual cycles using self-tracked mobile-health data,''
\textit{Nature Digital Medicine}, vol.~3, 2020,
arXiv:1808.02932.

\bibitem{isayev2023}
A. Isayev, N. Velieva, L. Isedisha, Z. Isayeva,
K. Kamburo\u{g}lu, and F. Kuyumcu,
``Cone-beam computed tomography as a prediction tool for osteoporosis
in postmenopausal women: A systematic literature review,''
\textit{Diagnostics}, vol.~13, p.~1027, 2023,
doi: 10.3390/diagnostics13061027.

\bibitem{rego2023}
R. Rego,
``Predictive modeling of menstrual cycle length:
A time series forecasting approach,''
arXiv preprint arXiv:2308.07927, 2023.

\bibitem{dataset}
``FemCare Reproductive Health Dataset for PCOS Risk Assessment
and Menstrual Cycle Prediction,''
Published September~25, 2026.
\url{https://zenodo.org/records/22953562}.

\end{thebibliography}
